% Options for packages loaded elsewhere
% Options for packages loaded elsewhere
\PassOptionsToPackage{unicode}{hyperref}
\PassOptionsToPackage{hyphens}{url}
\PassOptionsToPackage{dvipsnames,svgnames,x11names}{xcolor}
%
\documentclass[
  defaultstyle,
  11pt]{thesis}
\usepackage{xcolor}
\usepackage{amsmath,amssymb}
\setcounter{secnumdepth}{5}
\usepackage{iftex}
\ifPDFTeX
  \usepackage[T1]{fontenc}
  \usepackage[utf8]{inputenc}
  \usepackage{textcomp} % provide euro and other symbols
\else % if luatex or xetex
  \usepackage{unicode-math} % this also loads fontspec
  \defaultfontfeatures{Scale=MatchLowercase}
  \defaultfontfeatures[\rmfamily]{Ligatures=TeX,Scale=1}
\fi
\usepackage{lmodern}
\ifPDFTeX\else
  % xetex/luatex font selection
\fi
% Use upquote if available, for straight quotes in verbatim environments
\IfFileExists{upquote.sty}{\usepackage{upquote}}{}
\IfFileExists{microtype.sty}{% use microtype if available
  \usepackage[]{microtype}
  \UseMicrotypeSet[protrusion]{basicmath} % disable protrusion for tt fonts
}{}
\makeatletter
\@ifundefined{KOMAClassName}{% if non-KOMA class
  \IfFileExists{parskip.sty}{%
    \usepackage{parskip}
  }{% else
    \setlength{\parindent}{0pt}
    \setlength{\parskip}{6pt plus 2pt minus 1pt}}
}{% if KOMA class
  \KOMAoptions{parskip=half}}
\makeatother
% Make \paragraph and \subparagraph free-standing
\makeatletter
\ifx\paragraph\undefined\else
  \let\oldparagraph\paragraph
  \renewcommand{\paragraph}{
    \@ifstar
      \xxxParagraphStar
      \xxxParagraphNoStar
  }
  \newcommand{\xxxParagraphStar}[1]{\oldparagraph*{#1}\mbox{}}
  \newcommand{\xxxParagraphNoStar}[1]{\oldparagraph{#1}\mbox{}}
\fi
\ifx\subparagraph\undefined\else
  \let\oldsubparagraph\subparagraph
  \renewcommand{\subparagraph}{
    \@ifstar
      \xxxSubParagraphStar
      \xxxSubParagraphNoStar
  }
  \newcommand{\xxxSubParagraphStar}[1]{\oldsubparagraph*{#1}\mbox{}}
  \newcommand{\xxxSubParagraphNoStar}[1]{\oldsubparagraph{#1}\mbox{}}
\fi
\makeatother

\usepackage{color}
\usepackage{fancyvrb}
\newcommand{\VerbBar}{|}
\newcommand{\VERB}{\Verb[commandchars=\\\{\}]}
\DefineVerbatimEnvironment{Highlighting}{Verbatim}{commandchars=\\\{\}}
% Add ',fontsize=\small' for more characters per line
\usepackage{framed}
\definecolor{shadecolor}{RGB}{241,243,245}
\newenvironment{Shaded}{\begin{snugshade}}{\end{snugshade}}
\newcommand{\AlertTok}[1]{\textcolor[rgb]{0.68,0.00,0.00}{#1}}
\newcommand{\AnnotationTok}[1]{\textcolor[rgb]{0.37,0.37,0.37}{#1}}
\newcommand{\AttributeTok}[1]{\textcolor[rgb]{0.40,0.45,0.13}{#1}}
\newcommand{\BaseNTok}[1]{\textcolor[rgb]{0.68,0.00,0.00}{#1}}
\newcommand{\BuiltInTok}[1]{\textcolor[rgb]{0.00,0.23,0.31}{#1}}
\newcommand{\CharTok}[1]{\textcolor[rgb]{0.13,0.47,0.30}{#1}}
\newcommand{\CommentTok}[1]{\textcolor[rgb]{0.37,0.37,0.37}{#1}}
\newcommand{\CommentVarTok}[1]{\textcolor[rgb]{0.37,0.37,0.37}{\textit{#1}}}
\newcommand{\ConstantTok}[1]{\textcolor[rgb]{0.56,0.35,0.01}{#1}}
\newcommand{\ControlFlowTok}[1]{\textcolor[rgb]{0.00,0.23,0.31}{\textbf{#1}}}
\newcommand{\DataTypeTok}[1]{\textcolor[rgb]{0.68,0.00,0.00}{#1}}
\newcommand{\DecValTok}[1]{\textcolor[rgb]{0.68,0.00,0.00}{#1}}
\newcommand{\DocumentationTok}[1]{\textcolor[rgb]{0.37,0.37,0.37}{\textit{#1}}}
\newcommand{\ErrorTok}[1]{\textcolor[rgb]{0.68,0.00,0.00}{#1}}
\newcommand{\ExtensionTok}[1]{\textcolor[rgb]{0.00,0.23,0.31}{#1}}
\newcommand{\FloatTok}[1]{\textcolor[rgb]{0.68,0.00,0.00}{#1}}
\newcommand{\FunctionTok}[1]{\textcolor[rgb]{0.28,0.35,0.67}{#1}}
\newcommand{\ImportTok}[1]{\textcolor[rgb]{0.00,0.46,0.62}{#1}}
\newcommand{\InformationTok}[1]{\textcolor[rgb]{0.37,0.37,0.37}{#1}}
\newcommand{\KeywordTok}[1]{\textcolor[rgb]{0.00,0.23,0.31}{\textbf{#1}}}
\newcommand{\NormalTok}[1]{\textcolor[rgb]{0.00,0.23,0.31}{#1}}
\newcommand{\OperatorTok}[1]{\textcolor[rgb]{0.37,0.37,0.37}{#1}}
\newcommand{\OtherTok}[1]{\textcolor[rgb]{0.00,0.23,0.31}{#1}}
\newcommand{\PreprocessorTok}[1]{\textcolor[rgb]{0.68,0.00,0.00}{#1}}
\newcommand{\RegionMarkerTok}[1]{\textcolor[rgb]{0.00,0.23,0.31}{#1}}
\newcommand{\SpecialCharTok}[1]{\textcolor[rgb]{0.37,0.37,0.37}{#1}}
\newcommand{\SpecialStringTok}[1]{\textcolor[rgb]{0.13,0.47,0.30}{#1}}
\newcommand{\StringTok}[1]{\textcolor[rgb]{0.13,0.47,0.30}{#1}}
\newcommand{\VariableTok}[1]{\textcolor[rgb]{0.07,0.07,0.07}{#1}}
\newcommand{\VerbatimStringTok}[1]{\textcolor[rgb]{0.13,0.47,0.30}{#1}}
\newcommand{\WarningTok}[1]{\textcolor[rgb]{0.37,0.37,0.37}{\textit{#1}}}

\usepackage{longtable,booktabs,array}
\usepackage{calc} % for calculating minipage widths
% Correct order of tables after \paragraph or \subparagraph
\usepackage{etoolbox}
\makeatletter
\patchcmd\longtable{\par}{\if@noskipsec\mbox{}\fi\par}{}{}
\makeatother
% Allow footnotes in longtable head/foot
\IfFileExists{footnotehyper.sty}{\usepackage{footnotehyper}}{\usepackage{footnote}}
\makesavenoteenv{longtable}
\usepackage{graphicx}
\makeatletter
\newsavebox\pandoc@box
\newcommand*\pandocbounded[1]{% scales image to fit in text height/width
  \sbox\pandoc@box{#1}%
  \Gscale@div\@tempa{\textheight}{\dimexpr\ht\pandoc@box+\dp\pandoc@box\relax}%
  \Gscale@div\@tempb{\linewidth}{\wd\pandoc@box}%
  \ifdim\@tempb\p@<\@tempa\p@\let\@tempa\@tempb\fi% select the smaller of both
  \ifdim\@tempa\p@<\p@\scalebox{\@tempa}{\usebox\pandoc@box}%
  \else\usebox{\pandoc@box}%
  \fi%
}
% Set default figure placement to htbp
\def\fps@figure{htbp}
\makeatother





\setlength{\emergencystretch}{3em} % prevent overfull lines

\providecommand{\tightlist}{%
  \setlength{\itemsep}{0pt}\setlength{\parskip}{0pt}}



 


%% before-body.tex — Quarto template partial for thesis.cls integration
%%
%% PURPOSE:
%% The CU Boulder thesis.cls expects specific preamble commands (\author, \degree,
%% \dept, \advisor, etc.) that are normally placed in a .tex file BEFORE
%% \begin{document}. Quarto generates the \begin{document} automatically, so we
%% inject these commands via Quarto's "include-before-body" mechanism — but they
%% actually need to appear in the PREAMBLE. We use "include-in-header" for the
%% preamble commands and "include-before-body" for any post-\begin{document} content.
%%
%% HOW TO USE:
%% This file is referenced in _quarto-thesis.yml via:
%%   include-in-header: ../inst/thesis/thesis-preamble.tex
%%
%% CUSTOMIZATION:
%% Edit the values below to match YOUR dissertation. These map to the commands
%% that thesis.cls uses to generate the title page, approval page, and abstract.
%%
%% For other universities: replace these commands with whatever your thesis.cls
%% expects in the preamble. The file name and Quarto config reference stay the same.

%% ============================================================
%% DISSERTATION METADATA — Edit these for your dissertation
%% ============================================================

%% Your dissertation title (appears on title page and abstract)
\title{Longitudinal Achievement Patterns in Reading and Mathematics: An AI-Native Dissertation Framework}

%% Your name: \author{First M.}{Last}
\author{Your}{Name}

%% Other degrees you hold (each on its own line, separated by \\)
\otherdegrees{B.A., Your University, 20XX \\
              M.S., Another University, 20XX}

%% The degree you are receiving:
%%   First arg: full name of the degree
%%   Second arg: abbreviation and field
\degree{Doctor of Philosophy}
       {Ph.D., Education Research and Methodology}

%% Your department:
%%   First arg: organizational designation (e.g., "Department of", "School of", "College of")
%%   Second arg: name of the department
\dept{School of}
     {Education}

%% Your dissertation advisor:
%%   First arg: title (e.g., "Prof.", "Dr.", "Associate Professor")
%%   Second arg: full name
\advisor{Prof.}
        {Your Advisor}

%% Committee member who signs the thesis (required)
\reader{Committee Member Two}

%% Additional committee members (optional — comment out if not needed)
\readerThree{Committee Member Three}
\readerFour{Committee Member Four}
%\readerFive{Committee Member Five}

%% Optional: IRB protocol number (uncomment if applicable)
% \IRBprotocol{E927F29.001X}

%% ============================================================
%% ABSTRACT — Edit with your actual abstract
%% ============================================================

\abstract{
This dissertation examines longitudinal achievement patterns in Reading and
Mathematics using a multi-year, multi-district assessment dataset. Three
research questions guide the investigation: (1) What are the longitudinal
achievement trends and how do proficiency rates change over time?
(2) How does student achievement vary across districts? (3) How do cohorts
of students progress through grade levels?

Beyond the substantive findings, this work demonstrates an AI-native
dissertation framework --- a reproducible, forkable system that unifies
dissertation text, statistical analyses, and interactive dissemination
into a single repository. The framework treats AI as a collaborator in
analysis, writing, and dissemination, while maintaining full transparency
and human authorship of all scholarly conclusions.
}

%% ============================================================
%% DEDICATION and ACKNOWLEDGMENTS (optional)
%% ============================================================

\dedication[Dedication]{
  To all who believe in open, reproducible scholarship.
}

\acknowledgements{ \OnePageChapter
  The author gratefully acknowledges the support of advisors, colleagues,
  and the open-source community. This framework builds on the work of
  many contributors to R, Quarto, and the broader data science ecosystem.
}

%% ============================================================
%% TABLE OF CONTENTS / LIST OF FIGURES / LIST OF TABLES
%% ============================================================
%% Use the "Short" versions if the list fits on one page.
%% Use the "empty" versions if there are no figures/tables.

\ToCisShort        % Use for 1-page Table of Contents
% \LoFisShort      % Use for 1-page List of Figures
% \emptyLoF        % Use if there is no List of Figures
% \LoTisShort      % Use for 1-page List of Tables
% \emptyLoT        % Use if there is no List of Tables
%% thesis-quarto-fix.tex
%% =====================
%% Fixes conflicts between Quarto's auto-generated LaTeX and thesis.cls.
%%
%% PROBLEM: Quarto always emits \title{}, \author{}, \date{}, and
%% \maketitle just before \begin{document}. But thesis.cls expects
%% these to be set via its own preamble commands (in thesis-preamble.tex)
%% and generates title page, abstract, dedication, etc. internally
%% at \begin{document} time.
%%
%% SOLUTION: After thesis-preamble.tex sets all the correct metadata,
%% we save thesis.cls's internal state and restore it after Quarto
%% overwrites \title/\author. We also make \maketitle a no-op since
%% thesis.cls handles title page generation itself.

%% Save the thesis title set by thesis-preamble.tex before Quarto
%% overwrites it with the last chapter's title.
\makeatletter
\let\saved@thesis@title\@title
\let\saved@thesis@author\@author

%% At \begin{document}, restore the thesis.cls title/author and
%% suppress Quarto's \maketitle (thesis.cls handles this internally).
\AddToHook{begindocument/before}{%
  \global\let\@title\saved@thesis@title
  \global\let\@author\saved@thesis@author
  \renewcommand{\maketitle}{}%
}
\makeatother
\makeatletter
\@ifpackageloaded{caption}{}{\usepackage{caption}}
\AtBeginDocument{%
\ifdefined\contentsname
  \renewcommand*\contentsname{Table of contents}
\else
  \newcommand\contentsname{Table of contents}
\fi
\ifdefined\listfigurename
  \renewcommand*\listfigurename{List of Figures}
\else
  \newcommand\listfigurename{List of Figures}
\fi
\ifdefined\listtablename
  \renewcommand*\listtablename{List of Tables}
\else
  \newcommand\listtablename{List of Tables}
\fi
\ifdefined\figurename
  \renewcommand*\figurename{Figure}
\else
  \newcommand\figurename{Figure}
\fi
\ifdefined\tablename
  \renewcommand*\tablename{Table}
\else
  \newcommand\tablename{Table}
\fi
}
\@ifpackageloaded{float}{}{\usepackage{float}}
\floatstyle{ruled}
\@ifundefined{c@chapter}{\newfloat{codelisting}{h}{lop}}{\newfloat{codelisting}{h}{lop}[chapter]}
\floatname{codelisting}{Listing}
\newcommand*\listoflistings{\listof{codelisting}{List of Listings}}
\makeatother
\makeatletter
\makeatother
\makeatletter
\@ifpackageloaded{caption}{}{\usepackage{caption}}
\@ifpackageloaded{subcaption}{}{\usepackage{subcaption}}
\makeatother
\usepackage{bookmark}
\IfFileExists{xurl.sty}{\usepackage{xurl}}{} % add URL line breaks if available
\urlstyle{same}
\hypersetup{
  pdftitle={Discussion},
  colorlinks=true,
  linkcolor={blue},
  filecolor={Maroon},
  citecolor={Blue},
  urlcolor={Blue},
  pdfcreator={LaTeX via pandoc}}


\title{Discussion}
\author{}
\date{}
\begin{document}
\maketitle


\section{Introduction}\label{introduction}

\subsection{Background}\label{sec-background}

Educational assessment data provides the foundation for understanding
student achievement across schools, districts, and states. Standardized
assessments in reading and mathematics generate scale scores that enable
comparisons of student performance relative to grade-level expectations.
When examined longitudinally --- across multiple years and grade levels
--- these data reveal patterns of achievement that inform educational
policy, resource allocation, and instructional practice.

Achievement levels derived from scale scores (such as Beginning,
Developing, Proficient, and Advanced) provide an accessible framework
for categorizing student performance. Tracking the proportion of
students reaching proficiency over time, across content areas, and
between districts offers a practical lens on educational equity and
effectiveness.

\subsection{Problem Statement}\label{sec-problem}

Despite the widespread availability of assessment data, systematic
analyses that examine achievement patterns across multiple dimensions
--- time, geography, grade level, and content area --- remain
underutilized in educational research. This dissertation addresses this
gap using 368,301 longitudinal assessment records spanning 5 academic
years across 3 districts.

\subsection{Research Questions}\label{sec-research-questions}

This dissertation is guided by three research questions:

\begin{enumerate}
\def\labelenumi{\arabic{enumi}.}
\tightlist
\item
  \textbf{RQ1}: What are the longitudinal achievement trends in Reading
  and Mathematics, and how do proficiency rates change over time?
\item
  \textbf{RQ2}: How does student achievement vary across districts, and
  what patterns of between-district inequality are evident?
\item
  \textbf{RQ3}: How do cohorts of students progress through grade
  levels, and what do achievement trajectories reveal about the
  stability of student performance?
\end{enumerate}

\subsection{Significance}\label{sec-significance}

This work contributes to the field in three ways. First, it provides
empirical evidence on longitudinal achievement patterns using a
multi-year, multi-district assessment dataset. Second, it demonstrates
analytic approaches for examining between-district variation in student
outcomes. Third, the entire analysis is implemented as a reproducible,
AI-native framework that other researchers can fork, customize, and
extend for their own dissertation research.

\subsection{Organization of the Dissertation}\label{sec-organization}

The remainder of this dissertation is organized as follows. Chapter 2
reviews the literature on assessment-based accountability, achievement
measurement, and educational equity. Chapter 3 describes the data
sources, measures, and analytic methods. Chapter 4 presents the results
organized by research question. Chapter 5 discusses the findings,
implications, limitations, and directions for future research.

\section{Literature Review}\label{literature-review}

\subsection{Overview}\label{sec-lit-overview}

This chapter reviews the literature relevant to understanding student
achievement measurement and patterns in K-12 education. The review is
organized into four sections: the role of standardized assessment in
educational accountability, approaches to longitudinal achievement
analysis, equity and between-district variation, and the emerging
intersection of reproducible research practices with educational data
science.

\subsection{Standardized Assessment and
Accountability}\label{sec-assessment-accountability}

Standardized assessments in reading and mathematics serve as the
backbone of educational accountability systems across the United States.
These assessments generate scale scores that enable comparisons of
student performance across schools, districts, and states. Performance
levels --- typically ranging from below basic to advanced --- provide an
accessible framework for communicating achievement relative to
grade-level standards.

The use of proficiency rates as an accountability metric has been both
widespread and contested. While proficiency-based measures offer clear
thresholds for adequate performance, critics argue that they can obscure
meaningful variation among students above or below the cut point and can
incentivize instructional triage toward students near the proficiency
boundary.

\subsection{Longitudinal Achievement
Analysis}\label{sec-longitudinal-analysis}

Longitudinal analysis of assessment data --- tracking achievement
patterns across multiple years and grade levels --- provides richer
information than cross-sectional snapshots. By examining trends over
time, researchers can distinguish transient fluctuations from sustained
patterns and can identify whether improvements in achievement are
broad-based or concentrated in particular subgroups or regions.

Score distribution analysis using percentile breakpoints (e.g., the
10th, 25th, 50th, 75th, and 90th percentiles) provides a more complete
picture of achievement than mean scores alone. Changes in the
distribution shape over time can reveal whether gains are driven by
improvements at the top, middle, or bottom of the performance continuum.

\subsection{Equity and Between-District Variation}\label{sec-equity}

A growing body of literature examines disparities in educational
outcomes across schools and districts. Between-district variation in
achievement reflects differences in resources, instructional quality,
student demographics, and community characteristics. Examining how
achievement distributions differ across districts --- rather than
focusing solely on mean differences --- provides a more nuanced
understanding of educational inequality.

District-level comparisons of proficiency rates, mean scores, and score
distributions contribute to identifying both high-performing districts
whose practices merit study and struggling districts that may benefit
from targeted support.

\subsection{Summary and Gaps}\label{sec-lit-summary}

The literature reveals several gaps that this dissertation addresses.
First, while individual assessment results are widely reported,
systematic longitudinal analyses examining achievement trends across
multiple years, content areas, and organizational levels remain
relatively sparse. Second, between-district achievement variation is
often examined only through mean comparisons rather than full
distributional analysis. Third, the integration of reproducible research
practices with modern AI tools represents an emerging frontier in
educational research methodology, offering the potential to make complex
analyses more transparent and accessible.

\section{Methods}\label{methods}

\subsection{Overview}\label{sec-methods-overview}

This chapter describes the data, measures, and analytic methods used to
address the research questions. All analyses are implemented as
reproducible R functions within the \texttt{dissertation} R package,
ensuring complete transparency and replicability.

\subsection{Data Source}\label{sec-data-source}

The dataset comprises 368,301 longitudinal assessment records for 67,547
students across grades 10, 3, 4, 5, 6, 7, 8, 9 in MATHEMATICS and
READING spanning 5 academic years (2020\_2021, 2021\_2022, 2022\_2023,
2023\_2024, 2024\_2025).

\subsection{Sample Description}\label{sec-sample}

\begin{longtable}[]{@{}llrr@{}}

\caption{\label{tbl-sample}Sample Size by Content Area and Year}

\tabularnewline

\toprule\noalign{}
CONTENT\_AREA & YEAR & N\_Students & N\_Observations \\
\midrule\noalign{}
\endhead
\bottomrule\noalign{}
\endlastfoot
MATHEMATICS & 2020\_2021 & 35,620 & 35,620 \\
MATHEMATICS & 2021\_2022 & 36,458 & 36,458 \\
MATHEMATICS & 2022\_2023 & 37,041 & 37,041 \\
MATHEMATICS & 2023\_2024 & 37,640 & 37,640 \\
MATHEMATICS & 2024\_2025 & 37,965 & 37,965 \\
READING & 2020\_2021 & 35,404 & 35,404 \\
READING & 2021\_2022 & 36,229 & 36,229 \\
READING & 2022\_2023 & 36,807 & 36,807 \\
READING & 2023\_2024 & 37,411 & 37,411 \\
READING & 2024\_2025 & 37,726 & 37,726 \\

\end{longtable}

The sample includes students from 3 districts and 121 schools.
Table~\ref{tbl-sample} presents the sample composition by content area
and year.

\subsection{Measures}\label{sec-measures}

\subsubsection{Scale Scores}\label{scale-scores}

Student achievement is measured using scale scores in Reading and
Mathematics. Scale scores range from NA to NA across the dataset.

\subsubsection{Achievement Levels}\label{achievement-levels}

Scale scores are classified into performance levels using grade-specific
cut points. The levels in the dataset are: Advanced, No Score, Partially
Proficient, Proficient, Unsatisfactory. Students at the Proficient or
Advanced level are considered to have met grade-level expectations.

\subsubsection{Organizational Variables}\label{organizational-variables}

Records include district identifiers (\texttt{DISTRICT\_NUMBER}) and
school identifiers (\texttt{SCHOOL\_NUMBER}), enabling analyses of
between-district and between-school variation.

\subsection{Analytic Approach}\label{sec-analytic-approach}

The analyses are organized to address each research question:

\textbf{RQ1 --- Achievement Trends}: Longitudinal achievement patterns
are examined by computing mean scores and proficiency rates by year and
content area using the \texttt{analyzeAssessment("achievement\_trends")}
function.

\textbf{RQ2 --- District Variation}: Between-district comparisons use
\texttt{analyzeAssessment("district\_comparison")} to examine variation
in mean scores, proficiency rates, and sample sizes across the 3
districts in the dataset.

\textbf{RQ3 --- Cohort Progress}: Cohort tracking via
\texttt{analyzeAssessment("cohort\_progress")} examines how students
with longitudinal records progress through grade levels over time.

Additionally, score distribution analyses using
\texttt{analyzeAssessment("score\_distribution")} provide percentile
breakpoints (10th, 25th, 50th, 75th, 90th) to characterize the full
shape of the achievement distribution.

\subsection{Reproducibility}\label{sec-reproducibility}

All analyses are implemented in the \texttt{dissertation} R package and
can be replicated by installing the package and calling the analysis
functions directly:

\begin{Shaded}
\begin{Highlighting}[]
\FunctionTok{library}\NormalTok{(dissertation)}
\FunctionTok{summarizeAssessment}\NormalTok{(}\AttributeTok{content\_area =} \StringTok{"READING"}\NormalTok{, }\AttributeTok{group\_by =} \StringTok{"GRADE"}\NormalTok{)}
\FunctionTok{analyzeAssessment}\NormalTok{(}\StringTok{"achievement\_trends"}\NormalTok{, }\AttributeTok{content\_area =} \StringTok{"MATHEMATICS"}\NormalTok{)}
\end{Highlighting}
\end{Shaded}

The complete analysis code, data, and documentation are available via
the package's GitHub repository and REST API.

\section{Results}\label{results}

\subsection{Overview}\label{sec-results-overview}

This chapter presents the findings organized by research question. All
results are generated programmatically from the \texttt{dissertation} R
package analysis functions, ensuring that the narrative is always
synchronized with the underlying data.

\subsection{RQ1: Achievement Trends}\label{sec-rq1}

Achievement trends across 5 academic years reveal the following
patterns:

\begin{longtable}[]{@{}
  >{\raggedright\arraybackslash}p{(\linewidth - 12\tabcolsep) * \real{0.1266}}
  >{\raggedright\arraybackslash}p{(\linewidth - 12\tabcolsep) * \real{0.1646}}
  >{\raggedleft\arraybackslash}p{(\linewidth - 12\tabcolsep) * \real{0.1392}}
  >{\raggedleft\arraybackslash}p{(\linewidth - 12\tabcolsep) * \real{0.1646}}
  >{\raggedleft\arraybackslash}p{(\linewidth - 12\tabcolsep) * \real{0.1899}}
  >{\raggedleft\arraybackslash}p{(\linewidth - 12\tabcolsep) * \real{0.1392}}
  >{\raggedleft\arraybackslash}p{(\linewidth - 12\tabcolsep) * \real{0.0759}}@{}}

\caption{\label{tbl-achievement-trends}Mean Score and Proficiency Rate
by Year and Content Area}

\tabularnewline

\toprule\noalign{}
\begin{minipage}[b]{\linewidth}\raggedright
YEAR
\end{minipage} & \begin{minipage}[b]{\linewidth}\raggedright
CONTENT\_AREA
\end{minipage} & \begin{minipage}[b]{\linewidth}\raggedleft
mean\_score
\end{minipage} & \begin{minipage}[b]{\linewidth}\raggedleft
median\_score
\end{minipage} & \begin{minipage}[b]{\linewidth}\raggedleft
pct\_proficient
\end{minipage} & \begin{minipage}[b]{\linewidth}\raggedleft
n\_students
\end{minipage} & \begin{minipage}[b]{\linewidth}\raggedleft
n
\end{minipage} \\
\midrule\noalign{}
\endhead
\bottomrule\noalign{}
\endlastfoot
2020\_2021 & MATHEMATICS & 543.4 & 548 & 56.6 & 35620 & 35620 \\
2020\_2021 & READING & 636.1 & 642 & 72.7 & 35404 & 35404 \\
2021\_2022 & MATHEMATICS & 545.9 & 550 & 58.3 & 36458 & 36458 \\
2021\_2022 & READING & 636.8 & 642 & 74.2 & 36229 & 36229 \\
2022\_2023 & MATHEMATICS & 548.0 & 552 & 59.0 & 37041 & 37041 \\
2022\_2023 & READING & 636.9 & 643 & 73.5 & 36807 & 36807 \\
2023\_2024 & MATHEMATICS & 546.2 & 549 & 57.8 & 37640 & 37640 \\
2023\_2024 & READING & 636.6 & 641 & 72.8 & 37411 & 37411 \\
2024\_2025 & MATHEMATICS & 546.0 & 551 & 59.0 & 37965 & 37965 \\
2024\_2025 & READING & 636.5 & 643 & 74.5 & 37726 & 37726 \\

\end{longtable}

Achievement trends across 5 year(s) show relative stability in mean
scale scores. The overall mean score is 591.2 with a proficiency rate of
65.8\%.

\subsubsection{Reading Achievement
Patterns}\label{reading-achievement-patterns}

Achievement trends across 5 year(s) show relative stability in mean
scale scores. The overall mean score is 636.6 with a proficiency rate of
73.5\%.

\subsubsection{Mathematics Achievement
Patterns}\label{mathematics-achievement-patterns}

Achievement trends across 5 year(s) show an upward trend in mean scale
scores. The overall mean score is 545.9 with a proficiency rate of
58.1\%.

\subsection{RQ2: District Variation}\label{sec-rq2}

\subsubsection{Between-District Achievement
Comparisons}\label{between-district-achievement-comparisons}

\begin{longtable}[]{@{}
  >{\raggedleft\arraybackslash}p{(\linewidth - 14\tabcolsep) * \real{0.1684}}
  >{\raggedright\arraybackslash}p{(\linewidth - 14\tabcolsep) * \real{0.1368}}
  >{\raggedleft\arraybackslash}p{(\linewidth - 14\tabcolsep) * \real{0.1158}}
  >{\raggedleft\arraybackslash}p{(\linewidth - 14\tabcolsep) * \real{0.1368}}
  >{\raggedleft\arraybackslash}p{(\linewidth - 14\tabcolsep) * \real{0.1579}}
  >{\raggedleft\arraybackslash}p{(\linewidth - 14\tabcolsep) * \real{0.1158}}
  >{\raggedleft\arraybackslash}p{(\linewidth - 14\tabcolsep) * \real{0.1053}}
  >{\raggedleft\arraybackslash}p{(\linewidth - 14\tabcolsep) * \real{0.0632}}@{}}

\caption{\label{tbl-district}Achievement by District}

\tabularnewline

\toprule\noalign{}
\begin{minipage}[b]{\linewidth}\raggedleft
DISTRICT\_NUMBER
\end{minipage} & \begin{minipage}[b]{\linewidth}\raggedright
CONTENT\_AREA
\end{minipage} & \begin{minipage}[b]{\linewidth}\raggedleft
mean\_score
\end{minipage} & \begin{minipage}[b]{\linewidth}\raggedleft
median\_score
\end{minipage} & \begin{minipage}[b]{\linewidth}\raggedleft
pct\_proficient
\end{minipage} & \begin{minipage}[b]{\linewidth}\raggedleft
n\_students
\end{minipage} & \begin{minipage}[b]{\linewidth}\raggedleft
n\_schools
\end{minipage} & \begin{minipage}[b]{\linewidth}\raggedleft
n
\end{minipage} \\
\midrule\noalign{}
\endhead
\bottomrule\noalign{}
\endlastfoot
470 & MATHEMATICS & 542.8 & 548 & 57.5 & 25022 & 49 & 69465 \\
470 & READING & 632.8 & 640 & 71.0 & 24780 & 49 & 68428 \\
1040 & MATHEMATICS & 568.0 & 573 & 69.7 & 23968 & 32 & 63606 \\
1040 & READING & 655.3 & 661 & 85.2 & 23962 & 32 & 63554 \\
2690 & MATHEMATICS & 522.7 & 526 & 44.9 & 18480 & 42 & 51653 \\
2690 & READING & 618.4 & 622 & 62.5 & 18475 & 42 & 51595 \\

\end{longtable}

District comparison reveals mean scale scores ranging from 522.7 to
655.3 across 3 districts (gap of 132.6 points). Proficiency rates range
from 44.9\% to 85.2\%.

\subsubsection{Achievement Summaries by
Grade}\label{achievement-summaries-by-grade}

\begin{longtable}[]{@{}
  >{\raggedright\arraybackslash}p{(\linewidth - 12\tabcolsep) * \real{0.0845}}
  >{\raggedleft\arraybackslash}p{(\linewidth - 12\tabcolsep) * \real{0.1549}}
  >{\raggedleft\arraybackslash}p{(\linewidth - 12\tabcolsep) * \real{0.1831}}
  >{\raggedleft\arraybackslash}p{(\linewidth - 12\tabcolsep) * \real{0.1268}}
  >{\raggedleft\arraybackslash}p{(\linewidth - 12\tabcolsep) * \real{0.2113}}
  >{\raggedleft\arraybackslash}p{(\linewidth - 12\tabcolsep) * \real{0.1549}}
  >{\raggedleft\arraybackslash}p{(\linewidth - 12\tabcolsep) * \real{0.0845}}@{}}

\caption{\label{tbl-grade-summary}Achievement Summary by Grade}

\tabularnewline

\toprule\noalign{}
\begin{minipage}[b]{\linewidth}\raggedright
GRADE
\end{minipage} & \begin{minipage}[b]{\linewidth}\raggedleft
mean\_score
\end{minipage} & \begin{minipage}[b]{\linewidth}\raggedleft
median\_score
\end{minipage} & \begin{minipage}[b]{\linewidth}\raggedleft
sd\_score
\end{minipage} & \begin{minipage}[b]{\linewidth}\raggedleft
pct\_proficient
\end{minipage} & \begin{minipage}[b]{\linewidth}\raggedleft
n\_students
\end{minipage} & \begin{minipage}[b]{\linewidth}\raggedleft
n
\end{minipage} \\
\midrule\noalign{}
\endhead
\bottomrule\noalign{}
\endlastfoot
10 & 642.4 & 652 & 80.4 & 53.7 & 22526 & 45553 \\
3 & 524.0 & 531 & 89.3 & 78.9 & 22923 & 44688 \\
4 & 549.7 & 557 & 80.2 & 74.6 & 22775 & 45491 \\
5 & 576.9 & 583 & 80.9 & 73.5 & 22572 & 45153 \\
6 & 588.3 & 596 & 81.5 & 69.1 & 22659 & 45403 \\
7 & 604.3 & 610 & 78.0 & 62.0 & 22917 & 45950 \\
8 & 616.1 & 621 & 74.8 & 60.5 & 22957 & 46010 \\
9 & 623.0 & 633 & 76.9 & 55.6 & 24498 & 50053 \\

\end{longtable}

\subsection{RQ3: Cohort Progress}\label{sec-rq3}

\begin{longtable}[]{@{}
  >{\raggedright\arraybackslash}p{(\linewidth - 10\tabcolsep) * \real{0.1333}}
  >{\raggedright\arraybackslash}p{(\linewidth - 10\tabcolsep) * \real{0.1733}}
  >{\raggedleft\arraybackslash}p{(\linewidth - 10\tabcolsep) * \real{0.1467}}
  >{\raggedleft\arraybackslash}p{(\linewidth - 10\tabcolsep) * \real{0.2000}}
  >{\raggedleft\arraybackslash}p{(\linewidth - 10\tabcolsep) * \real{0.1467}}
  >{\raggedleft\arraybackslash}p{(\linewidth - 10\tabcolsep) * \real{0.2000}}@{}}

\caption{\label{tbl-cohort}Cohort Progress: Students with Longitudinal
Records}

\tabularnewline

\toprule\noalign{}
\begin{minipage}[b]{\linewidth}\raggedright
YEAR
\end{minipage} & \begin{minipage}[b]{\linewidth}\raggedright
CONTENT\_AREA
\end{minipage} & \begin{minipage}[b]{\linewidth}\raggedleft
mean\_score
\end{minipage} & \begin{minipage}[b]{\linewidth}\raggedleft
pct\_proficient
\end{minipage} & \begin{minipage}[b]{\linewidth}\raggedleft
n\_students
\end{minipage} & \begin{minipage}[b]{\linewidth}\raggedleft
n\_observations
\end{minipage} \\
\midrule\noalign{}
\endhead
\bottomrule\noalign{}
\endlastfoot
2020\_2021 & MATHEMATICS & 538.6 & 61.3 & 28713 & 28713 \\
2020\_2021 & READING & 630.3 & 74.0 & 28491 & 28491 \\
2021\_2022 & MATHEMATICS & 546.8 & 58.9 & 35104 & 35104 \\
2021\_2022 & READING & 637.5 & 74.6 & 34880 & 34880 \\
2022\_2023 & MATHEMATICS & 549.0 & 59.4 & 35733 & 35733 \\
2022\_2023 & READING & 637.5 & 73.7 & 35504 & 35504 \\
2023\_2024 & MATHEMATICS & 547.3 & 58.2 & 36248 & 36248 \\
2023\_2024 & READING & 637.5 & 73.2 & 36046 & 36046 \\
2024\_2025 & MATHEMATICS & 557.6 & 56.6 & 30095 & 30095 \\
2024\_2025 & READING & 646.0 & 73.6 & 30071 & 30071 \\

\end{longtable}

Cohort analysis tracks 48,578 students with 2+ years of data. Across the
cohort, the mean scale score is 592.7 with an overall proficiency rate
of 66.3\%.

\subsection{Score Distribution Properties}\label{sec-distribution}

\begin{longtable}[]{@{}llrrrrrrrr@{}}

\caption{\label{tbl-distribution}Score Distribution by Content Area and
Grade}

\tabularnewline

\toprule\noalign{}
CONTENT\_AREA & GRADE & p10 & p25 & p50 & p75 & p90 & mean & sd & n \\
\midrule\noalign{}
\endhead
\bottomrule\noalign{}
\endlastfoot
MATHEMATICS & 10 & 502 & 552 & 600 & 642 & 676 & 593.1 & 70.3 & 22785 \\
MATHEMATICS & 3 & 376 & 425 & 476 & 527 & 581 & 478.1 & 83.4 & 22860 \\
MATHEMATICS & 4 & 415 & 458 & 503 & 547 & 588 & 501.9 & 71.0 & 22790 \\
MATHEMATICS & 5 & 443 & 487 & 532 & 576 & 617 & 531.1 & 71.2 & 22576 \\
MATHEMATICS & 6 & 454 & 497 & 545 & 592 & 635 & 544.4 & 73.3 & 22704 \\
MATHEMATICS & 7 & 476 & 519 & 565 & 609 & 648 & 562.3 & 69.0 & 22965 \\
MATHEMATICS & 8 & 494 & 536 & 578 & 618 & 654 & 574.7 & 65.4 & 23006 \\
MATHEMATICS & 9 & 489 & 537 & 586 & 629 & 664 & 579.3 & 71.4 & 25038 \\
READING & 10 & 624 & 663 & 698 & 728 & 754 & 691.9 & 55.6 & 22768 \\
READING & 3 & 497 & 537 & 575 & 612 & 648 & 572.2 & 67.3 & 21828 \\
READING & 4 & 530 & 568 & 603 & 634 & 661 & 597.6 & 56.8 & 22701 \\
READING & 5 & 549 & 589 & 627 & 663 & 694 & 622.8 & 61.9 & 22577 \\
READING & 6 & 555 & 598 & 639 & 674 & 703 & 632.3 & 63.6 & 22699 \\
READING & 7 & 570 & 612 & 652 & 688 & 717 & 646.3 & 62.4 & 22985 \\
READING & 8 & 583 & 624 & 663 & 698 & 726 & 657.6 & 59.0 & 23004 \\
READING & 9 & 600 & 637 & 672 & 703 & 728 & 666.8 & 53.8 & 25015 \\

\end{longtable}

Scale score distributions show a median of 598 with an interquartile
range of 120 points, indicating substantial variation in student
achievement.

\subsection{Summary of Findings}\label{sec-results-summary}

The results reveal several key patterns across the three research
questions. Achievement trends show the trajectory of proficiency rates
and mean scores over the years examined. District comparisons reveal
variation in outcomes that warrants further investigation. The cohort
analysis demonstrates how student achievement trajectories develop over
the longitudinal span of the dataset, while score distribution analyses
reveal the full shape of the achievement continuum beyond mean-level
summaries.

\section{Discussion}\label{discussion}

\subsection{Summary of Findings}\label{sec-summary}

This dissertation examined longitudinal student achievement patterns in
Reading and Mathematics using a multi-year, multi-district assessment
dataset. Three research questions guided the investigation, each
yielding findings that contribute to our understanding of achievement
patterns in educational measurement.

\subsection{Interpretation of Results}\label{sec-interpretation}

\subsubsection{Achievement Trends (RQ1)}\label{achievement-trends-rq1}

The achievement trend analyses revealed patterns in mean scores and
proficiency rates across academic years. Examining trends by content
area provides insight into whether improvements (or declines) are
broad-based or concentrated in particular subjects. The consistency or
divergence of Reading and Mathematics trajectories over time informs
discussions about the relative effectiveness of curricular and
instructional investments.

\subsubsection{District Variation (RQ2)}\label{district-variation-rq2}

The between-district comparisons uncovered variation in achievement that
reflects underlying differences in resources, instructional quality, and
community characteristics. These findings align with the literature on
educational inequality and suggest that district-level analyses provide
a practical unit for identifying both exemplary practice and areas
needing support.

\subsubsection{Cohort Progress (RQ3)}\label{cohort-progress-rq3}

The cohort tracking analysis demonstrated how students progress through
grade levels over time. Examining the achievement trajectories of
students with longitudinal records provides insight into the stability
of student performance and the extent to which early achievement
predicts later outcomes.

\subsection{Implications}\label{sec-implications}

\subsubsection{For Policy}\label{for-policy}

The findings suggest that longitudinal achievement analyses provide
richer information than cross-sectional snapshots for understanding
educational effectiveness. Policymakers should consider incorporating
multi-year trend data and distributional analyses alongside
point-in-time proficiency rates to capture a more complete picture of
system performance.

\subsubsection{For Practice}\label{for-practice}

Educators can use achievement trend and distribution data to identify
grade levels or content areas where performance is changing and to
target instructional resources accordingly. Between-district comparisons
can help districts benchmark their performance and identify peers from
whom to learn.

\subsubsection{For Research Methodology}\label{for-research-methodology}

The AI-native dissertation framework demonstrates a new paradigm for
scholarly work in which the analysis code, documentation, and
dissemination are unified in a single reproducible system. This approach
enhances transparency, replicability, and accessibility of research.

\subsection{Limitations}\label{sec-limitations}

Several limitations should be considered when interpreting these
findings. First, the dataset used in this framework demonstration is
intended to illustrate methodological patterns and the AI-native
framework architecture rather than to support specific policy
conclusions. Second, the analyses presented here are descriptive and do
not account for school- or teacher-level effects that may confound
district-level comparisons. Third, without demographic covariates in the
dataset, equity analyses at the student level are limited to
organizational comparisons (district, school) rather than individual
student characteristics.

\subsection{Future Directions}\label{sec-future}

Future research could extend this work in several directions:
incorporating multilevel models to account for the nested structure of
educational data, adding demographic variables to enable equity-focused
disaggregation, integrating growth measures alongside achievement
status, and leveraging the AI-native framework to enable more
interactive and exploratory analyses.

\subsection{Conclusion}\label{sec-conclusion}

This dissertation contributes both substantive findings on student
achievement patterns and a methodological framework for conducting
AI-native research. By open-sourcing the analysis code, documentation,
and web application, this work aims to lower the barrier for future
researchers to conduct rigorous, transparent, and accessible
dissertation research.

\section*{References}\label{references}
\addcontentsline{toc}{section}{References}




\end{document}
